\section{sums over arbitrary sets}
\index{series!over arbitrary sets}%
\index{sums!over arbitrary sets}%
%%%%%%%%%%

\begin{lemma}\label{lemma:12734}
  Let $A$ be a countable set and let 
  $\alpha\colon \NN\to A$ and $\beta\colon \NN\to A$ 
  be two bijective functions.
  Let $f\colon A\to \RR$ (or 
  $f\colon A \to \CC$) be a function such that 
  \begin{equation}\label{eq:12734}
    \sum_{k=0}^{+\infty} \abs{f(\alpha(k))} < +\infty,  
  \end{equation}
  or let $f\colon A\to [0,+\infty)$ be any function. 
  Then we have
  \begin{equation}\label{eq:12735}
    \sum_{k=0}^{+\infty} f(\alpha(k)) =
    \sum_{j=0}^{+\infty} f(\beta(j)). 
  \end{equation}
\end{lemma}
%
\begin{proof}
Set $a_k=f(\alpha (k))$ and $\sigma = \alpha^{-1}\circ \beta$, we have 
$a_{\sigma(j)} = f(\beta (j))$.
Thus, if~\eqref{eq:12734} holds, we can 
apply theorem~\ref{th:convergenza_incondizionata}
to obtain~\eqref{eq:12735}.

If instead $f(x)\ge 0$, we obviously have $\abs{f(x)}=f(x)$,
thus if at least one of the two sums in~\eqref{eq:12735}
is finite, we apply theorem~\ref{th:convergenza_incondizionata}
to the series $\sum a_k$ and thus obtain the equality 
in~\eqref{eq:12735}.
If instead both series are divergent
(since they are series with non-negative terms, there are no other
possibilities!)
the equality~\eqref{eq:12735} still holds
since $+\infty = +\infty$.
\end{proof}

\begin{definition}[arbitrary sums]
Let $f\colon A\to [0,+\infty]$ be a function defined on a countable set 
$A$ (i.e., $\#A \le \#\NN$).
Then we define
  \[
    \sum_{x\in A}f(x)
    = \sum_{k=0}^{+\infty} f(\sigma(k)) 
  \]
where $\sigma\colon \NN\to A$ is any bijective function
(the result does not depend on $\sigma$ thanks to lemma~\ref{lemma:12734}).

If $f\colon A\to \CC$ or $f\colon A\to \RR$ (still with $A$ countable),
we say that $f$ is summable if $\sum_{x\in A}\abs{f(x)}<+\infty$. 
In such a case, we can define $\sum_{x\in A} f(x)$ as before,
since even in this case the sum does not depend
on the order of the addends.

If $A=\ENCLOSE{a_0, \dots, a_{n-1}}$ is a finite set, we can define 
\[
  \sum_{x\in A} f(x) = \sum_{k=0}^{n-1} f(a_k)
\]
and it is clear that the sum does not depend on the order of the addends.
\end{definition}

\begin{theorem}[summable families]
    Let $A = \displaystyle\bigcup_{n\in \NN} A_n$ be a union of disjoint sets:
  $A_n\cap A_m=\emptyset$ if $n\neq m$.
  Suppose that $A$ is countable (or finite).
  Let $f\colon A \to \RR$ be non-negative or $f\colon A \to \CC$ 
  a summable function (i.e., $\displaystyle\sum_{x\in A}\abs{f(x)}< +\infty$).
  Then 
  \[
    \sum_{x\in A}f(x) = \sum_{n\in \NN} \sum_{x\in A_n} f(x).  
  \]
\end{theorem}
%
\begin{proof}
Let $x_k$ be an enumeration of the elements of $A$ and define 
the double-index sequence:
\[
  a_{n,k} = \begin{cases}
    f(x_k) & \text{if $x_k\in A_n$}\\
    0 & \text{otherwise.}
  \end{cases}  
\]
For each $k\in \NN$, the following sum 
has only one non-zero addend:
\[
 \sum_{n=0}^{+\infty}  a_{n,k} = f(x_k)
\]
while for each $n\in \NN$, we have 
\[
 \sum_{k=0}^{+\infty} a_{n,k} = \sum_{x\in A_n} f(x)  
\]
since the non-zero terms of the sequence 
$a_{n,k}$ are equal to the value of $f$ on an enumeration 
of the elements of $A_n$.

Then we have
\[
 \sum_{k=0}^{+\infty}\sum_{n=0}^{+\infty} \abs{a_{n,k}}
 = \sum_{k=0}^{+\infty} {f(x_k)} = \sum_{x\in A} \abs{f(x_k)}.
\]
If this last sum is finite or if the terms are all non-negative,
we can then interchange the two sums thanks to theorem~\ref{th:scambio_somma}:
\[
  \sum_{x\in A} f(x) 
  = \sum_{k=0}^{+\infty} \sum_{n=0}^{+\infty} a_{n,k}
  = \sum_{n=0}^{+\infty} \sum_{k=0}^{+\infty} a_{n,k}
  = \sum_{n\in \NN} \sum_{x\in A_n}f(x).  
\]
\end{proof}

%%%